\subsection{ENVELOPES OF A FAMILY OF REAL-VALUED FUNCTIONS}

DEFINITION 4. Let \((f_{t})_{t\in I}\) be a family of real-valued functions defined on a set X. The real-valued function on X whose value at each point \(x\in X\) is \(\sup_{t\in I}(f_t(x))\) {[}resp. inf \((f_t(x))]\) is called the upper (resp. lower) envelope of the family \((f_t)\) , and is denoted by \(\sup_{t\in I}f_{t}\) or \(\sup_{t}f_{t}\) (resp. inf \(f_{t}\) or inf \(f_{t}\) ).

The upper envelope of the family \((f_{t})\) is thus the least upper bound of this family in the lattice \(\overline{R}\) of real-valued functions defined on X, and this justifies the notation \(\sup f_{t}\) .

Furthermore, if we endow \(\overline{\mathbf{R}}^{\mathbf{X}}\) with the topology which is the product of the topologies of its factors (all identical with \(\overline{R}\) ), we have the following proposition:

PROPOSITION 10. In the product space \(\overline{\mathbf{R}}^{\mathbf{X}}\) the upper envelope \(\sup f_{t}\) of a family of real-valued functions \((f_{t})_{t\in I}\) is the limit, with respect to the directed set \(\mathfrak{F}(\mathbf{I})\) of finite subsets of \(\mathbf{I}\) , of the mapping \(\mathrm{H}\to \sup_{t\in \mathbf{H}}f_t\) {[}which maps each finite subset \(\mathrm{H}\) of \(\mathbf{I}\) to the upper envelope of the finite subfamily \((f_{t})_{t\in H}]\) .

This follows immediately from Proposition 5 of no. 4 and from Chapter I, § 7, no. 6, Corollary 1 to Proposition 10.

We may therefore write

\[
\sup _ {t \in I} f _ {t} = \lim _ {H \in \mathfrak {F} (I)} \left(\sup _ {t \in H} f _ {t}\right).\tag{9}
\]

DEFINITION 5. A family \((f_t)_{i\in I}\) of real-valued functions defined on a set \(\mathbf{X}\) is said to be uniformly bounded above (resp. uniformly bounded below) in \(\mathbf{X}\) , if there exists a finite number \(a\) such that \(f_{t}(x) \leqslant a\) {[}resp. \(f_{t}(x) \geqslant a\) {]} for all \(x \in X\) and all \(t \in I\) . The family \((f_{t})\) is said to be uniformly bounded in \(X\) if it is uniformly bounded above and below in \(X\) .

Thus \((f_{t})\) is uniformly bounded above in \(\mathbf{X}\) if and only if the upper envelope of this family is bounded above in \(\mathbf{X}\) . \((f_{t})\) is uniformly bounded in \(\mathbf{X}\) if and only if the upper envelope of the family \((|f_{t}|)\) is bounded above in \(\mathbf{X}\) {[}i.e.~if and only if there is a finite real number \(a \geqslant 0\) such that \(|f_{t}(x)| \leqslant a\) for all \(x \in \mathbf{X}\) and all \(t \in I\) {]}.

